% ============================================================================
% SEC04.TEX - Section 2.4: Workflow Pengembangan Game
% ============================================================================

\section{Workflow Pengembangan Game di Unity}

\begin{learningoutcomes}
\item Memahami workflow dasar pengembangan game di Unity
\item Mengetahui urutan kerja yang efisien
\item Memahami konsep Scene dan Prefab
\item Mampu mengorganisir project dengan baik
\end{learningoutcomes}

\subsection{Workflow Dasar}

\begin{enumerate}
\item \textbf{Planning}: Rencanakan struktur scene dan assets
\item \textbf{Asset Creation}: Buat atau import assets
\item \textbf{Scene Setup}: Atur GameObject dalam scene
\item \textbf{Scripting}: Tulis script untuk gameplay
\item \textbf{Testing}: Test game di Play mode
\item \textbf{Iteration}: Perbaiki dan refine
\item \textbf{Build}: Build untuk target platform
\end{enumerate}

\subsection{Konsep Scene}

\begin{definisi}[Scene]
Scene adalah container yang berisi semua GameObject, lighting, dan settings untuk level atau bagian tertentu dari game \cite{UnityManual2024}.
\end{definisi}

\textbf{Best Practices untuk Scene}:
\begin{itemize}
\item Satu scene per level atau area
\item Organize dengan folder dalam Hierarchy
\item Gunakan empty GameObject sebagai organizer
\item Save scene secara teratur
\end{itemize}

\subsection{Konsep Prefab}

\begin{definisi}[Prefab]
Prefab adalah template GameObject yang dapat digunakan kembali dan diinstansiasi berkali-kali \cite{UnityManual2024}.
\end{definisi}

\textbf{Kapan Menggunakan Prefab}:
\begin{itemize}
\item GameObject yang digunakan berulang kali
\item GameObject dengan konfigurasi kompleks
\item GameObject yang perlu diupdate secara konsisten
\end{itemize}

\subsection{Organisasi Project}

\textbf{Struktur Folder yang Disarankan}:
\begin{itemize}
\item \textbf{Scenes/}: Semua scene files
\item \textbf{Scripts/}: Semua C\# scripts
\item \textbf{Art/}: Models, textures, sprites
\item \textbf{Audio/}: Music dan sound effects
\item \textbf{Prefabs/}: Semua prefab
\item \textbf{Materials/}: Material files
\item \textbf{Animations/}: Animation clips
\end{itemize}

\subsection{Version Control}

Unity bekerja baik dengan Git:
\begin{itemize}
\item Gunakan .gitignore untuk Unity
\item Jangan commit Library/ dan Temp/
\item Commit Assets/ dan ProjectSettings/
\end{itemize}

\begin{catatan}
Organisasi yang baik sejak awal akan menghemat waktu di kemudian hari. Investasikan waktu untuk setup struktur folder yang rapi.
\end{catatan}
